% !TEX root =../ECE343_Slides_Rogers.tex


%%%%%%%%%%%%%%%%%%%%%%%%%%%%%%%%%%%%%%%%%%%%%%%%%%%%%%
%%%%%%%%%%%%%%%%%%%%%%%%%%%%%%%%%%%%%%%%%%%%%%%%%%%%%%

\section{Schedule}
\renewcommand{\arraystretch}{1}
\begin{frame}{}
	\begin{textblock}{14}(0.5,1.75)
	\def\myfontsize{footnotesize}
	
\setcounter{mylnum}{1}
\begin{center}
\begin{\myfontsize}
\begin{tabular}{| c | l | c | l |}
\hline
\textbf{Lesson} 		  	& \textbf{Topics}             & \pbox{5cm}{\textbf{Reading}}	& \textbf{Notes}		\\ \hline\hline
\ifnum\lessonNum=\value{mylnum} \rowcolor{green!20!white}\fi
	\themylnum  
	& Course Overview  \& Introduction		
	& Ch.1 
	&  LO 1.1-1.2 
	\stepcounter{mylnum}\\ \hline \hline
\rowcolor{gray!20!white}\multicolumn{4}{| c |}{\textbf{Chapter 2: Electric and Magnetic Fields}}	\\ \hline\hline
\ifnum\lessonNum=\value{mylnum} \rowcolor{green!20!white}\fi
	\themylnum 	
	& Fields, Parameters, Field Quantities	
	& Ch. 2	
	& LO 2.1-2.2			
\stepcounter{mylnum}\\ \hline\hline
\rowcolor{gray!20!white}\multicolumn{4}{| c |}{\textbf{Chapter 3: Transmission Lines}}	 \\ \hline\hline	
\ifnum\lessonNum=\value{mylnum} \rowcolor{green!20!white}\fi
	\themylnum 	
	& Lumped Element Model				
	& 3.1-3.4 
	& LO 3.1		
\stepcounter{mylnum}\\ \hline
\ifnum\lessonNum=\value{mylnum} \rowcolor{green!20!white}\fi
	\themylnum 	
	& Telegrapher's Equations
	& 3.5-3.6
	&  LO 3.1	
\stepcounter{mylnum}\\ \hline
\ifnum\lessonNum=\value{mylnum} \rowcolor{green!20!white}\fi
	\themylnum 	
	&  Characteristic Impedance \& Wave Propagation	
	& 3.7-3.9		
	& LO 3.2						
\stepcounter{mylnum}\\ \hline
\ifnum\lessonNum=\value{mylnum} \rowcolor{green!20!white}\fi
	\themylnum 	
	& Coaxial,  Reflection Coefficient, VSWR	
	& 3.10,3.12-3.14	
	& LO 3.2-3.3		
\stepcounter{mylnum}\\ \hline 
\ifnum\lessonNum=\value{mylnum} \rowcolor{green!20!white}\fi
	\themylnum 	
	&  Terminated Transmission Lines	
	&  3.15,3.17-3.19
	& LO 3.4		
\stepcounter{mylnum}\\ \hline 
\ifnum\lessonNum=\value{mylnum} \rowcolor{green!20!white}\fi
	\themylnum	
	& Impedance Matching	
	& 3.21-3.23
	& LO 3.4 			
\stepcounter{mylnum}\\ \hline 
\ifnum\lessonNum=\value{mylnum} \rowcolor{green!20!white}\fi
	\themylnum  
	& Impedance Matching	
	& 3.21-3.23
	& LO 3.4		
\stepcounter{mylnum}\\ \hline 
\ifnum\lessonNum=\value{mylnum} \rowcolor{green!20!white}\fi
	\themylnum	
	& Smith Charts
	& Supplemental
	& LO 3.5
\stepcounter{mylnum}\\ \hline 
\ifnum\lessonNum=\value{mylnum} \rowcolor{green!20!white}\fi
	\themylnum	
	& Smith Charts
	& Supplemental 
	& LO 3.5					
\stepcounter{mylnum}\\ \hline 
\ifnum\lessonNum=\value{mylnum} \rowcolor{green!20!white}\fi
	\themylnum	
	& Smith Charts
	& Supplemental 
	& LO 3.5	
\stepcounter{mylnum}\\ \hline 
\ifnum\lessonNum=\value{mylnum} \rowcolor{green!20!white}\fi
	\themylnum	
	& Smith Charts
	& Supplemental 
	& LO 3.5	
\stepcounter{mylnum}\\ \hline 

\ifnum\lessonNum=\value{mylnum} \rowcolor{green!20!white}\fi
	\themylnum	
	& Transients		
	& Supplemental 
	& LO 3.6	 
\stepcounter{mylnum}\\ \hline \hline
\ifnum\lessonNum=\value{mylnum} \rowcolor{green!20!white}\fi
	\themylnum 	
	& Flex Day
	& 
	& 	
\stepcounter{mylnum}\\ \hline 
\ifnum\lessonNum=\value{mylnum} \rowcolor{green!20!white}\fi
	\themylnum 	
	& Reflection/Transmission Coefficient Lab
	& Lab Packet
	& LO 3.7		
\stepcounter{mylnum}\\ \hline 
\ifnum\lessonNum=\value{mylnum} \rowcolor{green!20!white}\fi
	\themylnum 	
	& Reflection/Transmission Coefficient Lab
	& Lab Packet
	& LO 3.7		
\stepcounter{mylnum}\\ \hline 
\rowcolor{gray!20!white}\multicolumn{4}{| c |}{\textbf{Chapter 4: Vector Analysis}} \\ \hline\hline
\ifnum\lessonNum=\value{mylnum} \rowcolor{green!20!white}\fi
	\themylnum	
	& Vector \& Coordinate System Review
	& Ch. 4
	&LO 4.1-4.4 \textbf{GR 1 due}	
\stepcounter{mylnum}\\ \hline 
\rowcolor{gray!20!white}\multicolumn{4}{| c |}{\textbf{Chapter 5: Electrostatics}}\\ \hline\hline
\ifnum\lessonNum=\value{mylnum} \rowcolor{green!20!white}\fi
	\themylnum	
	& Coulomb's Law
	& 5.1-5.3
	& LO 5.1
\stepcounter{mylnum}\\ \hline 
\ifnum\lessonNum=\value{mylnum} \rowcolor{green!20!white}\fi
	\themylnum	
	& Surface and Volume Charge			
	& 5.4
	& LO 5.1
\stepcounter{mylnum}\\ \hline 
\ifnum\lessonNum=\value{mylnum} \rowcolor{green!20!white}\fi
	\themylnum	
	& Gauss' Law (Integral Form)
	& 5.5-5.6
	& LO 5.2		 					
\stepcounter{mylnum}\\ \hline 
\ifnum\lessonNum=\value{mylnum} \rowcolor{green!20!white}\fi
	\themylnum	
	& Gauss' Law (Differential Form)		
	& 5.7-5.8
	& LO 5.2	 	
\stepcounter{mylnum}\\ \hline 
\ifnum\lessonNum=\value{mylnum} \rowcolor{green!20!white}\fi
	\themylnum	
	& Electric Potential	
	& 5.9-5.10, 5.15
	&  LO 5.3	
\stepcounter{mylnum}\\ \hline 
% <---- line 129 should be 20 lessons, based on 6 lines per lesson, starting at line 9 
\end{tabular}			
\end{\myfontsize}
\end{center}

	\end{textblock}
\end{frame}

%%%%%%%%%%%%%%%%%%%%%%%%%%%%%%%%%%%%%%%%%%%%%%%%%%%%%
%%%%%%%%%%%%%%%%%%%%%%%%%%%%%%%%%%%%%%%%%%%%%%%%%%%%%

\def\obja{I can describe the origin of the transmission line equation, understand the method of developing a solution, and interpret the solution.}
\def\objb{I can explain the physical meaning of and calculate the following quantities: propagation constant, phase constant, attenuation constant, characteristic impedance, VSWR, reflection coefficient, and transmission coefficient.}
\def\objc{I can use the transmission line equation to calculate the voltage/current at an arbitrary point on a specified transmission line.}
\def\objd{ I can use the transmission line equation to match a load to a source using a length of transmission line.}
\def\obje{ I can identify the various quantities represented on the Smith Chart.}
\def\objf{ I can use a Smith Chart to design a matching network for a given source and load, using series components, parallel components, shorted stubs, and open stubs.}
\def\objg{ I can plot the transient response of a load connected to a source through a transmission line. }
\def\objh{ I can measure the reflection coefficient across a range of frequencies and interpret the measurement using a network analyzer.}
\begin{frame}{Ch 3 Learning Objectives}
	\begin{textblock}{14}(0.5,3)
	\begin{scriptsize}
	\begin{enumerate}[{3}.1]
		\ifnum\objNum=1
			\item \alert{ \obja }
		\else 
			\item \obja
		\fi 
		\ifnum\objNum=23
			\item \alert{ \objb }
			\item \alert{ \objc }
		\else 
			\item \objb
			\item \objc
		\fi 
%		\ifnum\objNum=3
%			\item \alert{ \objc }
%		\else 
%			\item \objc
%		\fi 
		\ifnum\objNum=4
			\item \alert{ \objd }
		\else 
			\item \objd
		\fi 
		\ifnum\objNum=56
			\item \alert{ \obje }
		\else 
			\item \obje
		\fi 
		\ifnum\objNum=56
			\item \alert{ \objf }
		\else 
			\item \objf
		\fi 
		\ifnum\objNum=7
			\item \alert{ \objg }
		\else 
			\item \objg
		\fi 
		\ifnum\objNum=8
			\item \alert{ \objh }
		\else 
			\item \objh
		\fi 

	\end{enumerate}
	\end{scriptsize}
	\end{textblock}
\end{frame}

%%%%%%%%%%%%%%%%%%%%%%%%%%%%%%%%%%%%%%%%%%%%%%%%%%%%%
%%%%%%%%%%%%%%%%%%%%%%%%%%%%%%%%%%%%%%%%%%%%%%%%%%%%%
\section{Introduction}
\begin{frame}{What is a Transmission Line?}
	\begin{textblock}{14.5}(0.5,3)

		\begin{myblock}[0.9]{Definition}	
			
			A transmission line is a structure intended to transport electromagnetic signals or power.

		\end{myblock}
	
		\visible<2->{
		\begin{myblock}[0.9]{Purpose}
			Transmission lines are designed to support guided waves with controlled impedance, low loss, and a degree of immunity from EMI.
		\end{myblock}
		} % end visible 2
	\end{textblock}

\end{frame}

%%%%%%%%%%%%%%%%%%%%%%%%%%%%%%%%%%%%%%%%%%%%%%%%%%%%%
%%%%%%%%%%%%%%%%%%%%%%%%%%%%%%%%%%%%%%%%%%%%%%%%%%%

\section{Channels}
\begin{frame}{Communication Channels}
	\setbeamercovered{invisible}
	\hspace{-1.25cm}\begin{tikzpicture}
	\definecolor{myblue}{RGB}{187,224,227}
	\draw[fill=blue!10!white, draw=white] (0,0) rectangle ++(6.65,10);
	\end{tikzpicture}
	
	\begin{textblock}{7.5}(0.5,3.55) % left blue bar space
		\itm{
			\setlength\itemsep{1cm}
			\item[]<1> Wired Transmission Lines
			\item[]<2> Optical Fibers
			\item[]<3> Waveguides
			\item[]<4> Wireless
		} % end itm
	\end{textblock}
	
	\begin{textblock}{7}(8.5,3) % right white space
		
		% Transmission Lines/Wired
		\only<1>{
		\textbf{\underline{Advantages}}
		\itm{
			\item Reliable
			\item Easy to repair
			\item Relatively secure
		} % end itm

		\vspace{0.5cm}
		
		\textbf{\underline{Disadvantages}}
		\itm{
			\item Wire resistance
			\item Limited frequency bandwidth
		} % end itm
		\hspace{-0.3cm}\center\includegraphics[scale=0.07]{ECE343/underwater_line}
		} % end only 1

		% fiber
		\only<2>{
		\textbf{\underline{Advantages}}
		\itm{
			\item Large bandwidth
			\item Flexible
			\item Relatively inexpensive
			\item Low attenuation
		} % end itm
		
		\vspace{0.5cm}
		
		\textbf{\underline{Disadvantages}}
		\itm{
			\item Difficulty of installation
			\item Repairability 
		} % end itm
		\center\includegraphics[scale=0.1]{ECE343/fiber}
		} % end only 2
		
		% Waveguides
		\only<3>{
		\textbf{\underline{Advantages}}
		\itm{
			\item Low attenuation
			\item Precise			
		} % end itm

		\vspace{0.5cm}
		
		\textbf{\underline{Disadvantages}}
		\itm{
			\item Heavy/Rigid
			\item Limited bandwidth
			\item Difficult to repair
			\item Expensive
		} % end itm
		\center\includegraphics[scale=0.09]{ECE343/waveguide}
		} % end only 3
		
		% wireless
		\only<4>{
		\textbf{\underline{Advantages}}
		\itm{
			\item Very little infrastructure required
			
		} % end itm
		
		\vspace{0.5cm}
		\textbf{\underline{Disadvantages}}
		\itm{
			\item Line-of-sight considerations
			\item Interference
			\item Security
			\item Large power required
		} % end itm
		\center\includegraphics[scale=0.6]{ECE343/wireless_transmission}
		} % end only 4

	\end{textblock}
\end{frame}


%%%%%%%%%%%%%%%%%%%%%%%%%%%%%%%%%%%%%%%%%%%%%%%%%%%%%
%%%%%%%%%%%%%%%%%%%%%%%%%%%%%%%%%%%%%%%%%%%%%%%%%%%%%

\begin{frame}{Waveguides}
	\begin{textblock}{15}(0.5,3)
%		\center\includegraphics[scale=0.5]{ECE343/wg_cross_section}
		\begin{tikzpicture}[scale=1]
			\pic {annotated cuboid={width=20, height=10, depth=30, units=}};
			\node at (-1,-1.5) {$a$};
			\node at (-2.5,-0.5) {$b$};
			\node at (1,-0.5) {$l$};
		\begin{scope}[xshift=-2cm, yshift=1.5cm]
			\draw[line] (0,0,0) --++ (1,0,0);
			\draw[line] (0,0,0) --++ (0,1,0);
			\draw[line] (0,0,0) --++ (0,0,1);
			\node at (1.25,0,0) {\scriptsize $x$};
			\node at (0,1.25,0) {\scriptsize $y$};
			\node at (0,0,1.25) {\scriptsize $-z$};
		\end{scope}
		\end{tikzpicture}
		\includegraphics[scale=0.75]{ECE343/waveguide_modes_2}
		
		\vspace{0.5cm}
		
		\center cutoff frequency $\longrightarrow f_{c} = \dfrac{c}{2a} \longrightarrow$ HPF!!!
	\end{textblock}
\end{frame}

%%%%%%%%%%%%%%%%%%%%%%%%%%%%%%%%%%%%%%%%%%%%%%%%%%%%%
%%%%%%%%%%%%%%%%%%%%%%%%%%%%%%%%%%%%%%%%%%%%%%%%%%%%%
\begin{frame}{Why Consider Transmission Lines Now?}
	\begin{textblock}{14.5}(0.5,4.5)

	\begin{itemize}
		\setlength\itemsep{18pt}
		
		\item We use them often, treatment up to this point has been simplified
		
		\item When ratio of tx line length ($\ell$) to wavelength ($\lambda$)
		
		\item Wave solutions to transmission line equations
		
		\item Reflection, transmission, and interference concepts
	
	\end{itemize}
	
	\end{textblock}

\end{frame}

%%%%%%%%%%%%%%%%%%%%%%%%%%%%%%%%%%%%%%%%%%%%%%%%%%%%%
%%%%%%%%%%%%%%%%%%%%%%%%%%%%%%%%%%%%%%%%%%%%%%%%%%%%%
\begin{frame}{Transmission Line Model}
	\begin{textblock}{14.5}(0.5,3)

	\begin{circuitikz}
		
		\draw (0,0) to [vco=\myv{s}] ++(0,3)
					to[twoport,a=\myz{s}] ++(1.5,0)
					to[short, -o] ++(1,0)
					to[tline] ++(4,0)
					to[short,o-] ++(1.5,0)
					to[twoport, a=\myz{L}] ++(0,-3)
					to[short,-o] ++(-1.5,0)
					to[tline] ++(-3.85,0)
					to[short,o-] ++(-2.65,0);
		\draw[<->] (2.5,3.5) --++ (4,0) node[midway, above] {$\ell$};
		\node at (4,1.5) {\myz{0}};
	\end{circuitikz}
	
	\end{textblock}

\end{frame}

%%%%%%%%%%%%%%%%%%%%%%%%%%%%%%%%%%%%%%%%%%%%%%%%%%%%%
%%%%%%%%%%%%%%%%%%%%%%%%%%%%%%%%%%%%%%%%%%%%%%%%%%%%%

\begin{frame}{Lumped Element Model}
	\begin{textblock}{9}(0.5,3)

%	\begin{center}
	
	\scalebox{0.65}{
	\begin{circuitikz}
		
		\foreach\x in {0,2,4,...,8}{
			\draw (\x,0) to [short, o-o] ++(2,0);
			\draw (\x,3) to [short, o-o] ++(2,0);
			\draw[fill=white] ({\x+0.5},-0.5) rectangle ++(1,4);
			\draw[<->] ({\x-0.1},-0.75) --++(1.9,0) node[midway, below] {$\Delta z$};
			}
			
	\end{circuitikz}
	
	} % end scalebox
%	\end{center}
	\end{textblock}
	
	\begin{textblock}{9}(0.5,9)
%	\begin{center}
	\scalebox{0.75}{
	\begin{circuitikz}
	
		\draw (0,3) to[R=$R^{\prime}\Delta z$] ++(2,0)
				to[L=$L^{\prime}\Delta z$] ++(2,0)
				to[short] ++(6,0)
				to[open,o-o] ++(0,-3)
				to[short] ++(-10,0)
				to[open,o-o] ++(0,3);
		\draw (6,3) to[R=$G^{\prime}\Delta z$] ++(0,-3);
		
		\draw (8,3) to[C=$C^{\prime}\Delta z$] ++(0,-3);
		
		\draw[<->] (0,-0.5) --++ (10,0) node[midway, below] {$\Delta z$};
	\end{circuitikz}
	
		} % end scalebox
%	\end{center}
	\end{textblock}
	
	\only<2->{
	% the work on the side..
	\begin{textblock}{9}(10,3)
		\scalebox{0.8}{
		\begin{circuitikz}
			\coordinate (O) at (0,0);
			\draw[pattern=north west lines, pattern color=black!60!white] (O) circle (2cm);
			\draw[fill=white] (O) circle (1cm); 
			\draw[->] (O) -- (45:1cm) node[midway, below] {$a$};
			\draw[->] (O) -- (80:2cm) node[above=0.25cm] {$b$};
%			\draw[->] (O) -- (-45:1.5cm) node[pos=0.35, below] {$\rho$};
			
			\draw (180:1cm) to[short] ++ (-1.5cm,0)
							to[battery,v_=\myv{0}] ++ (0,-1cm)
							to[short] ++ (0.65cm,0);
			
			\draw (-2.25cm, -1cm) to[short] ++(0,-1cm) node[ground] {};
			\draw (-1cm, 1cm) to[short] ++ (-1cm,1cm)  node[above] {Dielectric};
		\end{circuitikz}
		} % end scalebox
		
	\end{textblock}
	} % end only 3
\end{frame}


%%%%%%%%%%%%%%%%%%%%%%%%%%%%%%%%%%%%%%%%%%%%%%%%%%%%%
%%%%%%%%%%%%%%%%%%%%%%%%%%%%%%%%%%%%%%%%%%%%%%%%%%%%%
\begin{frame}{Lumped Element Model: KVL}
	\begin{textblock}{14.5}(0.5,3)
%	\begin{center}
	\scalebox{0.75}{
	\begin{circuitikz}
		\draw (0,3) to[R=$R^{\prime}\Delta z$] ++(2,0)
				to[L=$L^{\prime}\Delta z$] ++(2,0)
				to[short] ++(6,0)
				to[open,o-o] ++(0,-3)
				to[short] ++(-10,0)
				to[open,o-o] ++(0,3);
		\draw (6,3) to[R=$G^{\prime}\Delta z$] ++(0,-3);
		
		\draw (8,3) to[C=$C^{\prime}\Delta z$] ++(0,-3);
		
		\draw[<->] (0,-0.5) --++ (10,0) node[midway, below] {$\Delta z$};
	\end{circuitikz}
		} % end scalebox
%	\end{center}
	\end{textblock}
	
	\begin{textblock}{14.5}(0.5,9)
	KVL:
	\vspace{-8pt}
	\alert{
	\eq{
		\visible<1->{-V(z,t) + R^{\prime} \Delta z i(z,t) &+ L^{\prime} \Delta z \partd{i(z,t)}{t} + V(z+\Delta z, t) = 0\\}
		 \visible<2->{\lim_{\Delta z \rightarrow 0} \dfrac{ V(z+\Delta z, t)-V(z,t)}{\Delta z} &= \dfrac{-R^{\prime} \Delta z i(z,t) - L^{\prime} \Delta z \partd{i(z,t)}{t} }{\Delta z}\\}
		 \visible<3->{-\partd{V(z,t)}{z} &= R^{\prime} i(z,t) + L^{\prime} \partd{i(z,t)}{t} \\}
	} % end eq
	} % end alert
	\end{textblock}
\end{frame}

%%%%%%%%%%%%%%%%%%%%%%%%%%%%%%%%%%%%%%%%%%%%%%%%%%%%%
%%%%%%%%%%%%%%%%%%%%%%%%%%%%%%%%%%%%%%%%%%%%%%%%%%%%%

\begin{frame}{Lumped Element Model: KCL}
	\begin{textblock}{14.5}(0.5,3)
%	\begin{center}
	\scalebox{0.75}{
	\begin{circuitikz}
		\draw (0,3) to[R=$R^{\prime}\Delta z$] ++(2,0)
				to[L=$L^{\prime}\Delta z$] ++(2,0)
				to[short] ++(6,0)
				to[open,o-o] ++(0,-3)
				to[short] ++(-10,0)
				to[open,o-o] ++(0,3);
		\draw (6,3) to[R=$G^{\prime}\Delta z$] ++(0,-3);
		
		\draw (8,3) to[C=$C^{\prime}\Delta z$] ++(0,-3);
		
		\draw[<->] (0,-0.5) --++ (10,0) node[midway, below] {$\Delta z$};
	\end{circuitikz}
		} % end scalebox
%	\end{center}
	\end{textblock}
	
	\begin{textblock}{14.5}(0.5,9)
	KVL:
	\vspace{-8pt}
	\alert{
	\eq{
		i(z,t) -i(z+\Delta z,t) - G^{\prime} \Delta z V(z+\Delta z,t)&-C^{\prime}\Delta z \partd{V(z+\Delta z, t)}{t}= 0\\
		&\vdots \\ 
	-\partd{i(z,t)}{z} &= G^{\prime} V(z,t) + C^{\prime} \partd{V(z,t)}{t} \\
	} % end eq
	} % end alert
	\end{textblock}
\end{frame}

%%%%%%%%%%%%%%%%%%%%%%%%%%%%%%%%%%%%%%%%%%%%%%%%%%%%%
%%%%%%%%%%%%%%%%%%%%%%%%%%%%%%%%%%%%%%%%%%%%%%%%%%%%%

\begin{frame}{Lumped Element Model}
	\begin{textblock}{14.5}(0.5,2.5)
		\begin{myblock}[0.9]{Telegrapher's Equations}
			\vspace{-8pt}
			\eq{
				-\partd{V(z,t)}{z} &= R^{\prime} i(z,t) + L^{\prime} \partd{i(z,t)}{t} \\
				-\partd{i(z,t)}{z} &= G^{\prime} V(z,t) + C^{\prime} \partd{V(z,t)}{t} 
			}
		\end{myblock}
	
	\only<2->{
	\begin{myblock}[0.9]{Phasor form of solutions}
	\vspace{-8pt}
			\eq{
				V(z,t) = \text{Re} \lb \widetilde{V}(z) e^{j\omega t} \rb 
				&\quad&
				i(z,t) = \text{Re} \lb \widetilde{I}(z) e^{j\omega t} \rb  \\
				\partd{V(z,t)}{t} = \partd{}{t} \text{Re} \lb \widetilde{V}(z) e^{j\omega t} \rb 
				&\quad&
				\partd{i(z,t)}{t} = \partd{}{t} \text{Re} \lb \widetilde{I}(z) e^{j\omega t} \rb \\
				= \text{Re} \lb j\omega \widetilde{V}(z) e^{j\omega t} \rb 
				&\quad& 
				= \text{Re} \lb j\omega \widetilde{I}(z) e^{j\omega t} \rb 
			}
		\end{myblock}
	} % end only 2
	\end{textblock}
\end{frame}

%%%%%%%%%%%%%%%%%%%%%%%%%%%%%%%%%%%%%%%%%%%%%%%%%%%%%
%%%%%%%%%%%%%%%%%%%%%%%%%%%%%%%%%%%%%%%%%%%%%%%%%%%%%

\begin{frame}{Decouple Telegrapher's Equations}
	\begin{textblock}{14.5}(0.5,3)
	Using Phasor form of solutions:
	\vspace{-12pt}
	\eqn{
		-\partd{\widetilde{V}(z)}{z} & = R^{\prime} \widetilde{I} (z) + j\omega L^{\prime} \widetilde{I}(z)  
		= \lb R^{\prime} + j\omega L^{\prime}\rb \widetilde{I}(z) \label{eq: t1}\\
			-\partd{\widetilde{I}(z)}{z} & = G^{\prime} \widetilde{V} (z) + j\omega C^{\prime} \widetilde{V}(z) 
			 = \lb G^{\prime} + j\omega C^{\prime}\rb \widetilde{V}(z) \label{eq: t2}
	} % end eqn 
	
	\vspace{-12pt}
	\visible<2->{Take $\partd{}{z}$ of (\ref{eq: t1})  to get:
	\vspace{-16pt}
		\eqn{
			-\partd{^2}{z^2} \widetilde{V}(z) &= \lb R^{\prime} + j\omega L^{\prime} \rb \partd{\widetilde{I}}{z} 
			\label{eq: t3}
		}
	} % end visible 2

	\vspace{-12pt}
	\visible<3->{
	Now, plug (\ref{eq: t2}) into (\ref{eq: t3}):
	\vspace{-12pt}
		\eqn{
			-\partd{^2}{z^2} \widetilde{V}(z) &= -\lb R^{\prime} + j\omega L^{\prime} \rb \lb G^{\prime} + j\omega C^{\prime} \rb \widetilde{V}(z)  \nn\\
			&\implies \partd{^2}{z^2} \widetilde{V}(z) - \lb R^{\prime} + j\omega L^{\prime} \rb \lb G^{\prime} + j\omega C^{\prime} \rb \widetilde{V}(z)  = 0  \nn\\
			&\implies \partd{^2}{z^2} \widetilde{V}(z) - \gamma^2 \widetilde{V}(z)  = 0 
			\text{  with: } \gamma = \sqrt{\lb R^{\prime} + j\omega L^{\prime} \rb \lb G^{\prime} + j\omega C^{\prime} \rb } \nn
		}
	} % end visible 3
	\end{textblock}
\end{frame}



%%%%%%%%%%%%%%%%%%%%%%%%%%%%%%%%%%%%%%%%%%%%%%%%%%%%%
%%%%%%%%%%%%%%%%%%%%%%%%%%%%%%%%%%%%%%%%%%%%%%%%%%%%%

\begin{frame}{Lumped Element Model: Wave Equations}
	\begin{textblock}{14.5}(0.5,3)
		\begin{myblock}[0.9]{Wave Equations}
			\vspace{-8pt}
			\eq{
				\partd{^2}{z^2} &\widetilde{V}(z) - \gamma^2 \widetilde{V}(z)  = 0 \\
				\partd{^2}{z^2} &\widetilde{I}(z) - \gamma^2 \widetilde{I}(z)  = 0  \\
				\gamma &=  \sqrt{\lb R^{\prime} + j\omega L^{\prime} \rb \lb G^{\prime} + j\omega C^{\prime} \rb}
			}
		\end{myblock}
	\end{textblock}
\end{frame}

%%%%%%%%%%%%%%%%%%%%%%%%%%%%%%%%%%%%%%%%%%%%%%%%%%%%%
%%%%%%%%%%%%%%%%%%%%%%%%%%%%%%%%%%%%%%%%%%%%%%%%%%%%%

\setbeamercovered{invisible}
\begin{frame}{Solving Wave Equations}
	\begin{textblock}{14.5}(0.5,2)
	\begin{myblock}[0.9]{Wave Equations}
			\vspace{-8pt}
			\eq{
				\partd{^2}{z^2} &\widetilde{V}(z) - \gamma^2 \widetilde{V}(z)  = 0 \\
				\partd{^2}{z^2} &\widetilde{I}(z) - \gamma^2 \widetilde{I}(z)  = 0  \\
				\gamma &=  \sqrt{\lb R^{\prime} + j\omega L^{\prime} \rb \lb G^{\prime} + j\omega C^{\prime} \rb}
				 = \alpha + j\beta 
			}
		\end{myblock}
		
		\begin{myblock}[0.9]{Solution Method}
			\begin{enumerate}[<+->]
%				\setlength\itemsep{12pt}
				\item Assume form of solution: \handoutgap{1in}{$A e^{\lambda z} $}
				\item Find characteristic equation \& it's \handoutgap{1in}{roots}
				\item Total solution is the sum of all possible solutions
				\item Enforce \handoutgap{1in}{boundary conditions} \&  \handoutgap{1in}{initial conditions}
			\end{enumerate}
		\end{myblock}
	\end{textblock}
\end{frame}

%%%%%%%%%%%%%%%%%%%%%%%%%%%%%%%%%%%%%%%%%%%%%%%%%%%%%
%%%%%%%%%%%%%%%%%%%%%%%%%%%%%%%%%%%%%%%%%%%%%%%%%%%%%

\begin{frame}{Wave Equation Solutions}
	\begin{textblock}{14.5}(0.5,2)
	\begin{myblock}[0.9]{Solution Forms}
			\vspace{-8pt}
			\eq{
				 \widetilde{V}(z)  &=\handoutgap{0in}{V^{+}_0 e^{-\gamma z} + V^{-}_0 e^{\gamma z}}\\[12pt]
				 \widetilde{I}(z)  &=  \handoutgap{0in}{I^{+}_0 e^{-\gamma z} +  I^{-}_0 e^{\gamma z}} \\[12pt]
				\gamma &=   \handoutgap{0in}{\sqrt{\lb R^{\prime} + j\omega L^{\prime} \rb \lb G^{\prime} + j\omega C^{\prime} \rb} } 
			}
		\end{myblock}
		

	\end{textblock}
\end{frame}


%%%%%%%%%%%%%%%%%%%%%%%%%%%%%%%%%%%%%%%%%%%%%%%%%%%%%
%%%%%%%%%%%%%%%%%%%%%%%%%%%%%%%%%%%%%%%%%%%%%%%%%%%%%

\begin{frame}{Back to Time Domain}
	\begin{textblock}{14.5}(0.5,3.5)
	To find the time domain solution, multiply by \handoutgap{0.5in}{$e^{j\omega t}$} and \handoutgap{0.5in}{take the real part}
	\visible<+->{
		\vspace{-8pt}
		\alert{
		\eq{
			V(z,t) &= \text{Re}\lb \widetilde{V}(z) e^{j\omega t} \rb = \text{Re} \lb V^{+}_0 e^{-\gamma z}e^{j\omega t} + V^{-}_0 e^{\gamma z}e^{j\omega t} \rb
		} % end eq
		} % end alert
	} % end visible 
	
	\vspace{-12pt}
	
	\visible<+->{
		Remember, $\gamma =$ \handoutgap{1in}{$\sqrt{\lb R^{\prime} + j\omega L^{\prime} \rb \lb G^{\prime} + j\omega C^{\prime} \rb} = \alpha + j\beta $}, which leads to:
	} % end visible

	\visible<+->{
		\vspace{-12pt}
		\alert{
		\eq{
			V(z,t) &= \text{Re} \lb V^{+}_0 e^{-\gamma z}e^{j\omega t} + V^{-}_0 e^{\gamma z}e^{j\omega t} \rb \\
				 &= \text{Re} \lb V^{+}_0 e^{-(\alpha + j\beta) z}e^{j\omega t} + V^{-}_0 e^{(\alpha + j\beta)z}e^{j\omega t} \rb \\				 
				&= \text{Re} \lb V^{+}_0 e^{-\alpha z}e^{j(\omega t -\beta z)}  V^{-}_0 e^{\alpha z}e^{j(\omega t +\beta z)} \rb
		} % end eq
		} % end alert
	} % end visible 

		\visible<+->{
		\vspace{-24pt}
		\begin{myblock}[0.9]{Instantaneous Solution to Wave Equations}
		\vspace{-12pt}
			\eq{
			V(z,t) &= V^{+}_0 e^{-\alpha z}\cosp{\omega t - \beta z} +  V^{-}_0 e^{\alpha z}\cosp{\omega t + \beta z}
			}
		\end{myblock}
		} % end visible
	\end{textblock}
\end{frame}

%%%%%%%%%%%%%%%%%%%%%%%%%%%%%%%%%%%%%%%%%%%%%%%%%%%%%
%%%%%%%%%%%%%%%%%%%%%%%%%%%%%%%%%%%%%%%%%%%%%%%%%%%%%

\begin{frame}{Interpretation of Solution}
	\begin{textblock}{8}(2.5,3.5)
		\vspace{-12pt}
			\eq{
				V(z,t) &=
				\underbrace{\alert<2>{V^{+}_0} \alert<3>{e^{-\alpha z}}\alert<4>{\cosp{\omega t - \beta z}}}_{\color{magenta}{\text{Forward Wave}}}
				+   \underbrace{\alert<2>{V^{-}_0} \alert<3>{e^{\alpha z}}\alert<4>{\cosp{\omega t + \beta z}}}_{\text{\color{green!50!black}{Reverse Wave}}}
				} % end eq 
			
	\end{textblock}

	\begin{textblock}{9}(0.5,6.5)
		\begin{itemize}
			\setlength\itemsep{12pt}
			\item<2-| alert@2> $V^{\pm}_{0} \longrightarrow$ Amplitude coefficients
			\item<3-| alert@3> $e^{\pm \alpha z} \longrightarrow$ Attenuation term
			\item<4-| alert@4> $\cosp{\omega t \pm \beta z} \longrightarrow$Time and space-varying wave
		\end{itemize}
	\end{textblock}
	
	\begin{textblock}{9}(8.5,6.5)
	\only<5>{
	\scalebox{0.75}{
	\begin{tikzpicture}
		\begin{axis}[domain=0:60e-3, 
			samples=200, 
			xmin=0, 
			xmax=65e-3, 
			ymax=4,
			ymin=-4,
	%		grid=major, 
			title={Forward Wave},
			scaled x ticks=base 10:3,
			xtick scale label code/.code={},
			minor x tick num = 1,
			minor y tick num = 1,
	%		ytick={-400,-300,-200,-100,0,100,200,300,400},
	%		xtick={-15e-3,-10e-3,-5e-3,0,5e-3,10e-3,15e-3,20e-3},
			yticklabel={\empty},
			xticklabel={\empty},
			xlabel=$z$,
			ylabel={$V(z,t)$},
	%		legend entries={$f=60$ Hz, $f=120$ Hz},
	%		legend pos=outer north east,
			axis y line = middle,
			axis x line = middle,
			x label style={xshift=0.5cm, yshift=-0.5cm},
			y label style={xshift=-0.5cm, yshift=0.75cm},
			clip=false,
			]
			\addplot [mark=none, orange!50!black, dashed] {cos(360*50*x)} 
				node[rounded corners,draw=orange!50!black,dashed,text=black,rectangle callout, fill=none, callout absolute pointer={(axis cs:12e-3,-1.1)}] at (axis cs: 20e-3, -2) {$\cosp{\omega t - \beta z}$};
			\addplot [dashed, magenta] {2*exp(-50*x)} 
				node[rounded corners,draw=magenta,dashed,text=black,rectangle callout, fill=none, callout absolute pointer={(axis cs:28e-3,0.5)}] at (axis cs: 25e-3, 2) {$e^{-\alpha z}$};
			\addplot[blue] {2*exp(-50*x)*cos(360*50*x)} 
				node[rounded corners,draw=blue,solid,text=black,rectangle callout, fill=none, callout absolute pointer={(axis cs:50e-3,-0.1)}] at (axis cs: 50e-3, 1.75) {$e^{-\alpha z}\cosp{\omega t - \beta z}$};

		\end{axis}
	\end{tikzpicture}
	} % end scalebox
	} % end only
	
	\only<6>{
	\scalebox{0.75}{
	\begin{tikzpicture}
		\begin{axis}[domain=0:60e-3, 
			samples=200, 
			xmin=0, 
			xmax=65e-3, 
			ymax=10,
			ymin=-10,
	%		grid=major, 
			title={Reverse Wave},
			scaled x ticks=base 10:3,
			xtick scale label code/.code={},
			minor x tick num = 1,
			minor y tick num = 1,
	%		ytick={-400,-300,-200,-100,0,100,200,300,400},
	%		xtick={-15e-3,-10e-3,-5e-3,0,5e-3,10e-3,15e-3,20e-3},
			yticklabel={\empty},
			xticklabel={\empty},
			xlabel=$z$,
			ylabel={$V(z,t)$},
	%		legend entries={$f=60$ Hz, $f=120$ Hz},
	%		legend pos=outer north east,
			axis y line = middle,
			axis x line = middle,
			x label style={xshift=0.5cm, yshift=-0.5cm},
			y label style={xshift=-0.5cm, yshift=0.75cm},
			clip=false,
			]
			\addplot [mark=none, yellow!50!black, dashed] {cos(360*50*x)};
				%node[rounded corners,draw=blue,solid,text=black,rectangle callout, fill=none, callout absolute pointer={(axis cs:12e-3,-1.1)}] at (axis cs: 20e-3, -2) {$\cosp{\omega t + \beta z}$};
			\addplot [dashed, magenta] {0.5*exp(50*x)} ;
%				node[rounded corners,draw=magenta,solid,text=black,rectangle callout, fill=none, callout absolute pointer={(axis cs:20e-3,1)}] at (axis cs: 25e-3, 2) {$e^{\alpha z}$};
			\addplot[blue] {0.5*exp(50*x)*cos(360*50*x)} ;
				%node[rounded corners,draw=orange!50!black,solid,text=black,rectangle callout, fill=none, callout absolute pointer={(axis cs:50e-3,-0.1)}] at (axis cs: 50e-3, 1.75) {$e^{\alpha z}\cosp{\omega t + \beta z}$};

		\end{axis}
	\end{tikzpicture}
	} % end scalebox
	} % end only
	
	\end{textblock}
	
	% do a graph
\end{frame}

%%%%%%%%%%%%%%%%%%%%%%%%%%%%%%%%%%%%%%%%%%%%%%%%%%%%%
%%%%%%%%%%%%%%%%%%%%%%%%%%%%%%%%%%%%%%%%%%%%%%%%%%%%%

\begin{frame}{Key Parameters of Waves}
	\begin{textblock}{8}(0.5,3)
		\scalebox{0.65}{
		\begin{tikzpicture}
			\begin{axis}[domain=0:60e-3, 
				samples=200, 
				xmin=0, 
				xmax=65e-3, 
				ymax=2,
				ymin=-2,
		%		grid=major, 
				title={$	V(z,t) = V_0^+ e^{-\alpha z}  \cosp{\omega t - \beta z}$},
				scaled x ticks=base 10:3,
				xtick scale label code/.code={},
				minor x tick num = 1,
				minor y tick num = 1,
		%		ytick={-400,-300,-200,-100,0,100,200,300,400},
		%		xtick={-15e-3,-10e-3,-5e-3,0,5e-3,10e-3,15e-3,20e-3},
				yticklabel={\empty},
				xticklabel={\empty},
				xlabel=$z$,
%				ylabel={$V(z,t)$},
		%		legend entries={$f=60$ Hz, $f=120$ Hz},
		%		legend pos=outer north east,
				axis y line = middle,
				axis x line = middle,
				x label style={xshift=0.5cm, yshift=-0.5cm},
				y label style={xshift=-0.5cm, yshift=0.75cm},
				clip=false,
				]
%				\addplot [mark=none, blue] {cos(360*50*x)} ;
%				\addplot [dashed, magenta] {2*exp(-50*x)} ;
				\addplot[blue] {2*exp(-50*x)*cos(360*50*x)} ;
%				\node at (axis cs: 35e-3,2) {$	V(z,t) = V_0^+ e^{-\alpha z}  \cosp{\omega t - \beta z}$};
			\end{axis}
		\end{tikzpicture}
		} % end scalebox
	\end{textblock}
	
	\begin{textblock}{8}(8,3)
	
	\begin{itemize}
	
		\ifdefined\ishandout
		 
		\else
			 \setlength\itemsep{60pt} 
		\fi
		
		\item Wavelength \handoutgap{1in}{$\lambda = \dfrac{2\pi}{\beta}$}
	
		\only<0|handout:1>{
		One period is $2\pi$ radians, so at $t=0$:
		\vspace{-12pt}
			\eq{
				\omega (0)- \beta z &= 2\pi \\
				z = \dfrac{2\pi}{\beta} = \lambda 
			} % end eq
		} % end only  
		
		\item Speed = \handoutgap{1in}{$\dfrac{\omega}{\beta}$}
	
		\only<0|handout:1>{
		Pick a point and ride the wave:
		\vspace{-12pt}
			\eq{
				\omega t - \beta z &= \text{constant} \\
				\omega t = \beta z \implies z &= \dfrac{\omega t}{\beta} \\
				u_p = v_p = \partd{z}{t} &= \dfrac{\omega}{\beta} \\[-24pt]
			} % end eq
		} % end only 
	\end{itemize}
	
	\end{textblock}
	
	\begin{textblock}{8.5}(0.5,10.25)
	
		\begin{itemize}
			
			\item Direction = \handoutgap{1in}{Forward}
			
			\only<0|handout:1>{
			Pick a point and ride the wave:
			\vspace{-12pt}
				\eq{
					\omega t - \beta z &= \text{constant} \\
					\omega t = \beta z \\[-24pt]
				} % end eq
			As $t$ increases, so does $z$
			} % end only 
			
		\end{itemize}
	
	\end{textblock}
\end{frame}

%%%%%%%%%%%%%%%%%%%%%%%%%%%%%%%%%%%%%%%%%%%%%%%%%%%%%
%%%%%%%%%%%%%%%%%%%%%%%%%%%%%%%%%%%%%%%%%%%%%%%%%%%%%

\begin{frame}{Key Parameters for Lossless Case}
	\begin{textblock}{14.5}(0.5,2)
	
	
	\begin{myblock}[0.9]{Lossless/low-loss Conditions}
		\vspace{-12pt}
		\eq{
			&\gamma = \sqrt{\lp R^{\prime} + j\omega C^{\prime}\rp\lp G^{\prime} + j\omega L^{\prime} \rp} \\
			&R^{\prime} \longrightarrow G^{\prime} \text{ or: } R^{\prime}<< \omega C^{\prime} \text{ and } 
			G^{\prime} << \omega L^{\prime}\\
			&\implies \gamma = \handoutgap{1.5in}{j\omega \sqrt{L^{\prime}C^{\prime}}} \\
			&\implies \alpha = \handoutgap{1in}{0} \quad \beta = \handoutgap{1in}{\omega\sqrt{L^{\prime}C^{\prime}}}
		}
	\end{myblock}
	
	\begin{myblock}[0.9]{Coax Cable}
		\vspace{-12pt}
		\eq{
			C^{\prime} &= \dfrac{2\pi\varepsilon}{\text{ln}\lp \frac{b}{a}\rp} 
			\qquad 
			L^{\prime} = \dfrac{\mu_0 }{2\pi} \text{ln}\lp \frac{b}{a}\rp \\
			&v_p = \handoutgap{1.5in}{\dfrac{1}{\sqrt{\varepsilon \mu}}}
		}
	\end{myblock}
	\end{textblock}
\end{frame}

%%%%%%%%%%%%%%%%%%%%%%%%%%%%%%%%%%%%%%%%%%%%%
%%%%%%%%%%%%%%%%%%%%%%%%%%%%%%%%%%%%%%%%%%%%%

\begin{frame}{Summary/Review}
	\begin{textblock}{14.5}(0.5,2)
	
	
	\begin{myblock}[0.9]{}
		\vspace{-12pt}
		\begin{scriptsize}
		\eq{
			\widetilde{V}(z) &= V^{+}_0 e^{-\gamma z}\ +  V^{-}_0 e^{\gamma z} \\
			\widetilde{I}(z) &= I^{+}_0 e^{-\gamma z} +  I^{-}_0 e^{\gamma z} \\
			-\partd{\widetilde{V}(z)}{z} & = \lp R^{\prime} + j\omega L^{\prime} \rp \widetilde{I}(z)  \\
			-\partd{\widetilde{I}(z)}{z} & = \lp G^{\prime} + j\omega C^{\prime} \rp \widetilde{V}(z)  
		}
		\end{scriptsize}
	\end{myblock}
	\vspace{-12pt}
	\begin{myblock}[1]{Characteristic Impedance - Forward Wave}
		\vspace{-12pt}
		\begin{scriptsize}
		\eq{
			\visible<1->{-\partd{\widetilde{V}(z)}{z} & =\lp  R^{\prime}+ j\omega L^{\prime} \rp \widetilde{I}(z) \\} 
			\visible<2->{ -\partd{}{z} \lp V^{+}_0 e^{-\gamma z}\ +  V^{-}_0 e^{\gamma z} \rp
				&=  \lp R^{\prime} + j\omega L^{\prime}  \rp \lp I^{+}_0 e^{-\gamma z} +  I^{-}_0 e^{\gamma z} \rp\\}
			\visible<3->{ \gamma V_0^+  &=  \lp R^{\prime}  + j\omega L^{\prime}  \rp I_0^+ \\ }
			\visible<4->{Z_0 &=  \dfrac{V_0^+}{I_0^+} =\sqrt{ \dfrac{ R^{\prime}  + j\omega L^{\prime}  }{\gamma} = \dfrac{R^{\prime}  + j\omega L^{\prime}}{G^{\prime}  + j\omega C^{\prime}}  }  }
		} % end eq
				\vspace{-12pt}
		\end{scriptsize}
	\end{myblock}
	\end{textblock}
\end{frame}


